% TOP_PROBLEM: {"id": 30006950, "title": "Capacity Region of General 3-Receiver Discrete Memoryless Broadcast Channels", "category_id": 15, "category": {"id": 15, "name": "computer_science", "display_name": "Computer Science", "description": "Computational complexity, algorithms, and theoretical CS.", "slug": "computer-science", "order_index": 15, "created_at": "2026-07-31T15:26:25.671Z"}, "status": "open"}
% FIRST_POSTED: 2026-09-25
% ENTRY_KIND: statement_only
% !TeX program = lualatex
% Definition and sources only; this file is not an LLM solution attempt.
\documentclass[11pt]{article}
\usepackage[margin=25mm]{geometry}
\usepackage{fontspec}
\setmainfont{Latin Modern Roman}
\usepackage{amsmath,amssymb,mathtools}
\usepackage{unicode-math}
\setmathfont{Latin Modern Math}
\usepackage{xurl,hyperref}
\hypersetup{colorlinks=true,urlcolor=blue}
\setcounter{secnumdepth}{0}
\setlength{\parindent}{0pt}
\setlength{\parskip}{5pt}
\setlength{\emergencystretch}{3em}
\begin{document}
\section*{\texorpdfstring{Capacity Region of General 3-Receiver Discrete Memoryless Broadcast Channels}{Capacity Region of General 3-Receiver Discrete Memoryless Broadcast Channels}}\label{30006950}
\subsection{Definitions and mathematical statement}
A three-receiver discrete memoryless broadcast channel has finite alphabets and transition law $W(y_1,y_2,y_3\mid x)$. Over $n$ uses, conditional output probabilities factor as the product of these single-use laws. An encoder maps independent messages $(M_1,M_2,M_3)$ to $X^n$, and receiver $i$ estimates $M_i$ from $Y_i^n$. A rate triple is achievable if a sequence of codes has vanishing average probability of decoding error and $n^{-1}\log_2|\mathcal M_i|$ approaching the corresponding rates.

Find a computable single-letter description of the closure of all achievable triples for every such channel, without degradedness assumptions. Single-letter means a finite-dimensional characterization in terms of one channel use and auxiliary variables with effective cardinality bounds, rather than a limit of arbitrary-block coding problems. The standard independent-private-message interpretation is used; adding common messages would require an additional rate specification not given in the record.
\par\smallskip\textit{Formulation and concept references:} \href{https://doi.org/10.1017/CBO9780511978166}{[S1]}.

\subsection{Short English statement}
What computable one-use information formula describes all achievable private-message rates for an arbitrary three-receiver broadcast channel?
\subsection{Sources}
\begin{itemize}
\item[S1]A. El Gamal, Y.-H. Kim, Network Information Theory, Cambridge University Press, 2011.
\url{https://doi.org/10.1017/CBO9780511978166}.
\textit{Formulation source.}
\end{itemize}
\end{document}
